\section{三角换元在椭圆中的应用}

\begin{definition}
    三角换元求值域，本质是在使用三角函数的有界性来求解问题。

    1．已知 $x^2+y^2=r^2$ ，求 $(m x+n y)_{\text {max}}$ ．联想到 $\sin ^2 \theta+\cos ^2 \theta=1$ ，通过观察不难设出 $x=r \cos \theta $、$y=r \sin \theta$ ，易得 $m x+n y=\sqrt{m^2 r^2+n^2 r^2} \sin (\theta+\varphi)$ ，即可求得答案。

    2．已知 $\left(\frac{x}{a}\right)^2+\left(\frac{y}{b}\right)^2=1$ ，求 $(m x+n y)_{\text {max }}$ ，同上可设 $x=a \cos \theta$ 、 $y=b \sin \theta$ 。
\end{definition}

\begin{example}
    点 $P(x, y)$ 在椭圆 $\frac{x^2}{4}+y^2=1$ 上移动，求 $u = 2x + 3y$ 的最大值和最小值.
\end{example}

\begin{tips}
    如果要求 $u=x^2+2 x y+4 y^2+x+2 y$ 的最大值呢？
\end{tips}

\begin{example}
    求椭圆 $\frac{x^2}{4}+y^2=1$ 上点 $P$ 到点 $Q(1, 2)$ 的距离的最大值和最小值.
\end{example}

\begin{tips}
    可以尝试用三角换元做一下。回忆：

    $$
        y=a \sin x \pm b \cos x=\sqrt{a^2+b^2} \sin (x \pm \varphi)
    $$

    其中 $\varphi$ 为锐角， $\tan \varphi=\frac{b}{a}$

    为了避免一些符号上的错误，因此我们约定：
    1． $\sin$ 在前； $2 . \sin$ 的系数为正； $3 . \varphi$ 为锐角．

\end{tips}